\section{Puzzles and scoring}\label{sec:scoring}

\subsection{Where did our puzzles come from?}

Previously aired puzzles in past episodes of Wheel of Fortune are on fan sites and are thus almost certainly in the models' training data. Therefore, we generate our puzzles that no model could've seen. We use MiniMax M3 (MiniMax) to write them and GLM-5.3 Flash (Zhipu AI) to judge them for category fit, familiarity, and naturalness, and if the category is a song, quote, or title, whether it's real (and not hallucinated). For the sake of fairness, these models do not play this game. 

Before that, though, generated puzzles must:
\begin{itemize}
  \item fit the board's four rows without splitting a word across rows,
  \item meet a minimum length,
  \item spell everything correctly as the bot's wordlists know it,
  \item meet basic category requirements,
  \item not be a duplicate or near duplicate of a previous puzzle, and
  \item if it's a bonus puzzle, reveal no more than 55\% of it after R S T L N E and the common picks of C, D, M, and A. 
\end{itemize}

Appendix~D lists the prompts and checks.

\subsection{The puzzle set}

The puzzle set that we used contains 400 main round puzzles and 100 bonus puzzles. To ensure representativeness of the puzzle set in accordance with how the show aired its puzzles across rounds and categories, we created our deck such that each puzzle is tagged with a round number and category. For example, when it comes to categories, ``What Are You Doing?'' puzzles always have an -ING word, ``Same Letter'' puzzles are always alliterations, and ``Rhyme Time'' puzzles contain two different words that rhyme. Our code checks and the judge model ensure compliance with category requirements like these.

In addition, we created the puzzle set such that the length and category distributions closely resemble that of previously aired puzzles. Our set conforms to median lengths of each of the 4 rounds and the bonus of 23, 22, 21, 17, and 13 letters respectively, against an archive of 22,933 previously aired puzzles from 2013 to 2025. Separate tuning decks were created for the bot during its development, which means the bot did not have an unfair advantage when it saw this set. Appendix~\ref{app:categories} contains more information about the puzzle distribution. 

The game engine as described in Section~\ref{sec:game} deals each round from the puzzles tagged with that round.


\subsection{Lineups and duplicate scoring}

A lineup is a ``tournament'' of 3 players, typically consisting of one LLM/decision model/hybrid against two copies of the bot. Each lineup plays \SetsPerLineup{} sets of 6 games each on seed \Seed{}. The 6 games correspond to each of the 6 possible permutations with these 3 seats. As the seed is shared, all 6 permutations meet the exact set of puzzles and wheel spin outcomes in the same order with only the seat order differing. Across different lineups, the shared seed makes each set directly comparable (i.e. the 40th set of lineup A can be compared to the 40th set of lineup B). 

This is known as duplicate scoring, a technique also used in bridge tournaments, to remove most of the luck of the deal from comparisons.

Even though two copies of the bot reduce unique permutations from 6 to 3 due to duplication, we still play as if there are 6 permutations as nondeterministic LLMs and decision models frequently occupy a seat, meaning that these models may play differently when asked to run it back. Therefore, the same seating arrangement is played twice to account for it. In total, we feature \NumLineups{} lineups that played \NumGames{} games in total.

\subsection{Measures}

The player with the highest total after the main rounds advances to the bonus round, and for all practical purposes, wins the game (as the bonus round is only played by the winner and can only add to the winner's total). In each lineup, we assume a ``null hypothesis'' where no player has an advantage when they each win 1/3 of the games. We draw 95\% confidence intervals from \Resamples{} resamples of these sets with fixed seeds. A model beats a bot when the 95\% interval of the set-by-set gap between them excludes zero. Section~\ref{sec:how-well} describes how these models did.
